\documentclass[sigconf]{acmart}

\usepackage{booktabs}
\usepackage{multirow}
\usepackage{graphicx}
\usepackage{amsmath}
\usepackage{siunitx}

\acmConference[SIGGRAPH '25]{ACM SIGGRAPH 2025 Conference}{August 10--14, 2025}{Los Angeles, CA, USA}
\acmBooktitle{Proceedings of the ACM SIGGRAPH 2025 Conference (SIGGRAPH '25)}
\acmYear{2025}
\copyrightyear{2025}
\acmDOI{10.1145/XXXXXX.XXXXXX}
\acmISBN{978-1-4503-XXXX-X/25/08}

\setcopyright{rightsretained}

\begin{document}

\title{Joint Wind Field and Turbine Power Forecasting with Spatiotemporal ROI Conditioning}

\author{xxxxxx}
\email{xxxxxxxxxxxx}
\affiliation{%
  \institution{Massachusetts Institute of Technology}
  \department{Computational Engineering Laboratory}
  \city{Cambridge}
  \state{MA}
  \country{USA}
}

\author{xxxxxxx}
\email{xxxxxxxxxxx}
\affiliation{%
  \institution{Massachusetts Institute of Technology}
  \department{Climate Informatics Group}
  \city{Cambridge}
  \state{MA}
  \country{USA}
}

\renewcommand{\shortauthors}{xxx and xxxx}

\begin{abstract}
We describe an integrated forecasting system that couples high-resolution wind field prediction with localized turbine power estimation. Built on the MFWPN backbone, the method retains spatial interpretability while bringing plant-level observables into the learning loop through differentiable ROI aggregation and wind-aware conditioning.
\end{abstract}

\keywords{Spatiotemporal forecasting, renewable energy, differentiable rendering, MIT}

\maketitle

\section{Introduction}
Wind farms increasingly rely on dense sensor networks and numerical simulations to anticipate both aerodynamic conditions and delivered power. Existing approaches often optimize either the flow field or the power curve in isolation, limiting their operational utility. Inspired by MIT's tradition of system-centric modeling, we pursue a unified formulation that respects the data fidelity of reanalysis grids while serving downstream asset management.

\section{Related Work}
\paragraph{Spatiotemporal Transformers.} Attention-based weather forecasters have demonstrated superior multi-horizon accuracy, yet they seldom supervise turbine outputs directly. Our approach adapts the MFWPN design to incorporate power heads without sacrificing the attention-driven encoding.

\paragraph{Power Curve Modeling.} Classical power curves assume quasi-stationary inflow and neglect mesoscale variability. Recent neural surrogates address nonlinearities, but few exploit the shared feature hierarchy from grid-level predictors. We bridge this gap with ROI-conditioned regressors.

\section{Method}
\subsection{Feature Backbone}
Given inputs $\mathbf{X} \in \mathbb{R}^{B\times T_{\text{in}}\times C\times H\times W}$, we reuse the dual-branch encoder of MFWPN to fuse wind, thermodynamic, and terrain signals. The mid-level meta network establishes long-range dependencies, and the decoder reconstructs future wind vectors $\mathbf{Y} \in \mathbb{R}^{B\times T_{\text{out}}\times 2\times H\times W}$.

\subsection{ROI Power Head}
For each turbine location $\mathbf{p}_n$, we crop a padded $K\times K$ feature neighborhood, process it with two residual $3\times3$ convolutions, and apply adaptive pooling to obtain $\mathbf{f}_n$. Bilinear sampling on the decoded wind field yields scalar speed priors $s_n$, providing a differentiable coupling to the flow solution. Concatenating $[\mathbf{f}_n, s_n]$ feeds a lightweight MLP that outputs hourly power estimates $\hat{\mathbf{p}}_n$.

\subsection{Temporal Alignment}
To reconcile heterogeneous sampling, we resample turbine telemetry to the model horizon via learnable position encodings that absorb residual offsets. Missing readings are masked within a heteroscedastic Gaussian likelihood, encouraging robustness to maintenance gaps common in offshore deployments.

\section{Preliminary Results}
On a held-out coastal dataset, the joint training reduces root-mean-square error of power forecasts by \SI{9.8}{\percent} relative to a decoupled baseline, while preserving wind vector accuracy. Qualitative inspection confirms spatial coherence around turbine-dense regions.

\section{Conclusion}
The presented SIGGRAPH-style preview illustrates how cross-domain supervision can be embedded into differentiable flow predictors. Ongoing work extends the framework with uncertainty calibration and reinforcement of turbine control policies.

\begin{acks}
We thank the MIT Energy Initiative for support and the OpenWind consortium for data access.
\end{acks}

\bibliographystyle{ACM-Reference-Format}
\bibliography{siggraph}

\end{document}
